\subsection{nml\_testcase\_vortex}
\label{SEC:nml_testcase_vortex}
\inp*[nmlTestcaseVortex]{}[nml\_testcase\_vortex][]

This \hyperref[SEC:User_Input]{namelist block}
is used to specify inputs for the vortex test cases.
%
These cases include the different versions of
  the isentropic Euler vortex problem
    ($\inp[nmlInitialization::Initialization]{Initialization}=[\pm5,\pm55,\pm6]$),
  the density profile advection case
    ($\inptext{Initialization}=7$),
  the entropy profile advection case
    ($\inptext{Initialization}=8$),
and
  the Taylor-Green vortex problem
    ($\inp{Initialization}=9$).

\begin{nmlblocklst}[
,label={LST:nml_testcase_vortex}
,caption={Namelist nml\_testcase\_vortex.}
]
&nml_testcase_vortex
  Vortex_Bigr          = 1.0   ! [m]
  Vortex_Biggam        = 5.0   ! [unitless]
  Grid_Lref            = 10.0  ! [m]
  Grid_Random_Perturb  = F
  Grid_Perturb_Percent = 30
  Grid_Hex_To_Pyramids = F
  Grid_Is_3d           = F
/
\end{nmlblocklst}

%==========================================================================

\begin{nmldescitem}
Vortex_Bigr = 1.0  ! [m]
\end{nmldescitem}
\inp*[nmlTestcaseVortex::VortexBigr]{}[Vortex\_Bigr][]

\begin{nmldescbody}
Constant used to specify the size or characteristic length (in meters)
of the isentropic Euler vortex.
\end{nmldescbody}

%==========================================================================

\begin{nmldescitem}
Vortex_Biggam = 5.0  ! [unitless]
\end{nmldescitem}
\inp*[nmlTestcaseVortex::VortexBiggam]{}[Vortex\_Biggam][]

\begin{nmldescbody}
Constant used to specify the strength of the isentropic Euler vortex.
\end{nmldescbody}

%==========================================================================

\begin{nmldescitem}
Grid_Lref = 10.0  ! [m]
\end{nmldescitem}
\inp*[nmlTestcaseVortex::GridLref]{}[Grid\_Lref][]

\begin{nmldescbody}
Reference length (in meters) for the grid domain,
\thatis the grid is a square with bounds of
$\left[ -\inptext{Grid\_Lref},
        +\inptext{Grid\_Lref}
 \right]$
in both the $x$- and $y$-directions.
\end{nmldescbody}

%==========================================================================

\begin{nmldescitem}
Grid_Random_Perturb = F
\end{nmldescitem}
\inp*[nmlTestcaseVortex::GridRandomPerturb]{}[Grid\_Random\_Perturb][]

\begin{nmldescbody}
Logical flag for applying a random perturbation to all the interior grid
nodes for the isentropic Euler vortex and Taylor-Green vortex problems.
\end{nmldescbody}

%==========================================================================

%##%\FINALFIX{\newpage}
\begin{nmldescitem}
Grid_Perturb_Percent = 30
\end{nmldescitem}
\inp*[nmlTestcaseVortex::GridPerturbPercent]{}[Grid\_Perturb\_Percent][]

\begin{nmldescbody}
\raggedright % prevent right margin overflow of input variable
Percent of the uniform grid cell length to perturb the interior grid nodes if
\inp[nmlTestcaseVortex::GridRandomPerturb]{Grid\_Random\_Perturb}
is true.
\end{nmldescbody}

%==========================================================================

\begin{nmldescitem}
Grid_Hex_To_Pyramids = F
\end{nmldescitem}
\inp*[nmlTestcaseVortex::GridHexToPyramids]{}[Grid\_Hex\_To\_Pyramids][]

\begin{nmldescbody}
When using the internal grid generator with one of the 3D vortex problems,
the default is to use hex cells. Setting this to true will create a grid
using only pyramid cells, by first creating the hex-only grid, and then
subdividing each hex cell into 6 pyramids by adding a new node point in the
center of the cell.
\end{nmldescbody}

%==========================================================================

\FINALFIX{\newpage}
\begin{nmldescitem}
Grid_Is_3d = F
\end{nmldescitem}
\inp*[nmlTestcaseVortex::GridIs3d]{}[Grid\_Is\_3d][]

\begin{nmldescbody}
When running one of the traditionally 2D vortex cases (\forexample
isentropic vortex), setting this to true will run a 3D version of the case
by extruding the domain and initial solution in $z$ to create a cube domain.
\end{nmldescbody}

%==========================================================================

